% ============================================================
%  Musaic — ERCİYES ÜNİVERSİTESİ Bitirme Ödevi (Final Raporu)
%  ERÜ Bilgisayar Mühendisliği resmi şablonuna uygun
% ============================================================
\documentclass[12pt, a4paper]{report}

% ─── Encoding ───────────────────────────────────────────────
\usepackage[utf8]{inputenc}
\usepackage[T1]{fontenc}
\usepackage[turkish, shorthands=off]{babel}

% ─── Layout ─────────────────────────────────────────────────
\usepackage[top=3cm, bottom=3cm, left=3.5cm, right=2.5cm]{geometry}

% ─── Typography ─────────────────────────────────────────────
\usepackage{microtype}
\usepackage{setspace}
\doublespacing

% ─── Math ───────────────────────────────────────────────────
\usepackage{amsmath, amssymb}

% ─── Tables & Figures ───────────────────────────────────────
\usepackage{graphicx}
\usepackage{booktabs}
\usepackage{tabularx}
\usepackage{multirow}
\usepackage{array}
\usepackage{float}
\usepackage{caption}
\usepackage{xcolor}
\usepackage{longtable}

% ─── TikZ ───────────────────────────────────────────────────
\usepackage{tikz}
\usetikzlibrary{shapes.geometric, arrows.meta, positioning, fit, backgrounds, shadows}

% ─── Listings ───────────────────────────────────────────────
\usepackage{listings}

% ─── Hyperref ───────────────────────────────────────────────
\usepackage[hidelinks]{hyperref}

% ─── Misc ───────────────────────────────────────────────────
\usepackage{fancyhdr}
\usepackage{enumitem}
\usepackage{titlesec}
\usepackage{tocloft}

% İçindekiler tablosunda BÖLÜM isimleri ile numaralar arasına nokta ekler
\renewcommand{\cftchapleader}{\cftdotfill{\cftdotsep}}

% ═══════════════════════════════════════════════════════════
%  RENK TANIMI
% ═══════════════════════════════════════════════════════════
\definecolor{primary}{RGB}{26, 60, 110}
\definecolor{accent}{RGB}{0, 140, 140}
\definecolor{lightbg}{RGB}{245, 248, 252}
\definecolor{codebg}{RGB}{250, 250, 252}

% ═══════════════════════════════════════════════════════════
%  BÖLÜM BAŞLIK STİLİ
%  Şablon formatı: "1. BÖLÜM" üstte, başlık altında — ortalı
% ═══════════════════════════════════════════════════════════
\titleformat{\chapter}[display]
  {\normalfont\large\bfseries\centering}
  {\large\bfseries\thechapter. BÖLÜM}
  {0pt}
  {\large\bfseries\MakeUppercase}

\titlespacing{\chapter}{0pt}{0pt}{24pt}

\titleformat{\section}
  {\normalfont\normalsize\bfseries}
  {\thesection.}{0.5em}{}

\titleformat{\subsection}
  {\normalfont\normalsize\bfseries}
  {\thesubsection.}{0.5em}{}

\titleformat{\subsubsection}
  {\normalfont\normalsize\bfseries}
  {\thesubsubsection.}{0.5em}{}

% ═══════════════════════════════════════════════════════════
%  HEADER / FOOTER — Sayfa numarası alt orta
% ═══════════════════════════════════════════════════════════
\pagestyle{fancy}
\fancyhf{}
\fancyfoot[C]{\thepage}
\renewcommand{\headrulewidth}{0pt}
\fancypagestyle{plain}{
  \fancyhf{}
  \fancyfoot[C]{\thepage}
  \renewcommand{\headrulewidth}{0pt}
}

% ═══════════════════════════════════════════════════════════
%  LISTINGS
% ═══════════════════════════════════════════════════════════
\lstset{
  backgroundcolor = \color{codebg},
  basicstyle      = \ttfamily\footnotesize,
  keywordstyle    = \bfseries\color{primary},
  commentstyle    = \itshape\color{gray},
  stringstyle     = \color{accent},
  breaklines      = true,
  frame           = single,
  framerule       = 0.4pt,
  rulecolor       = \color{gray!40},
  numbers         = left,
  numberstyle     = \tiny\color{gray},
  xleftmargin     = 1.8em,
  tabsize         = 2,
  showstringspaces= false,
}

% ─── Numaralandırma bölüm bazlı ─────────────────────────────
\renewcommand{\thetable}{\thechapter.\arabic{table}}
\renewcommand{\thefigure}{\thechapter.\arabic{figure}}

% ─── Caption etiketleri ─────────────────────────────────────
\renewcommand{\figurename}{Şekil}
\renewcommand{\tablename}{Tablo}

% ═══════════════════════════════════════════════════════════
\begin{document}
% ═══════════════════════════════════════════════════════════

% ────────────────────────────────────────────────────────────
%  KAPAK SAYFASI
% ────────────────────────────────────────────────────────────
\begin{titlepage}
  \thispagestyle{empty}
  \begin{center}
    {\large\bfseries T.C.}\\[0.3cm]
    {\large\bfseries ERCİYES ÜNİVERSİTESİ}\\[0.2cm]
    {\large\bfseries MÜHENDİSLİK FAKÜLTESİ}\\[3cm]

    {\LARGE\bfseries MUSAİC: YAPAY ZEKA DESTEKLİ DUYGU DUYARLI MÜZİK ÖNERİ SİSTEMİ}\\[3cm]

    {\normalsize
    \begin{tabular}{ll}
      \textbf{Hazırlayan} & \\[0.2cm]
      Berke Balcı         & 1030510273 \\[0.1cm]
      İlayda Baran        & 1030510572 \\[0.8cm]
      \textbf{Danışman}  & \\[0.2cm]
      Ömür Şahin          & \\
    \end{tabular}
    }

    \vfill
    {\normalsize\bfseries Bilgisayar Mühendisliği}\\[0.2cm]
    {\normalsize\bfseries Bitirme Ödevi}\\[0.5cm]
    {\normalsize Haziran 2026}\\[0.3cm]
    {\normalsize\bfseries KAYSERİ}
  \end{center}
\end{titlepage}

% ────────────────────────────────────────────────────────────
%  ÖN SAYFALAR — Romen rakamı numaralandırma
% ────────────────────────────────────────────────────────────
\pagenumbering{roman}
\setcounter{page}{1}

% ── KABUL VE ONAY ───────────────────────────────────────────
\newpage
\thispagestyle{plain}
\phantomsection
\addcontentsline{toc}{chapter}{KABUL VE ONAY}
\begin{center}
  \vspace*{1cm}
  \textit{``Musaic: Yapay Zeka Destekli Duygu Duyarlı Müzik Öneri Sistemi'' adlı bu çalışma, jürimiz tarafından Erciyes Üniversitesi Bilgisayar Mühendisliği Bölümünde Bitirme Ödevi olarak kabul edilmiştir.}

  \vspace{1.5cm}
  \hspace*{4cm}19/06/2026

  \vspace{1.5cm}
  \begin{flushleft}
    \textbf{JÜRİ :}

    \vspace{0.8cm}
    Danışman~: Ömür Şahin \hfill \underline{\hspace{5cm}}

    \vspace{0.8cm}
    Üye \hspace{1.1cm}: \hfill \underline{\hspace{5cm}}

    \vspace{0.8cm}
    Üye \hspace{1.1cm}: \hfill \underline{\hspace{5cm}}
  \end{flushleft}

  \vspace{2cm}
  \textbf{ONAY :}

  \vspace{0.5cm}
  Yukarıdaki imzaların, adı geçen öğretim elemanlarına ait olduğunu onaylıyorum.

  \vspace{3cm}
  \hfill .... / .... /2026

  \vspace{1cm}
  \hfill\begin{tabular}{c}
    Prof. Dr. Veysel ASLANTAŞ \\
    Bilgisayar Müh. Bölüm \\
    Başkanı
  \end{tabular}
\end{center}

% ── ÖNSÖZ / TEŞEKKÜR ────────────────────────────────────────
\newpage
\phantomsection
\addcontentsline{toc}{chapter}{ÖNSÖZ / TEŞEKKÜR}
\chapter*{ÖNSÖZ / TEŞEKKÜR}

Bu çalışma, Erciyes Üniversitesi Mühendislik Fakültesi Bilgisayar Mühendisliği Bölümü bitirme ödevi kapsamında hazırlanmıştır. Projenin fikir aşamasından, mimari tasarımına ve son kodlama süreçlerine kadar değerli görüşleri, teknik vizyonu ve yapıcı geri bildirimleriyle bize yol gösteren kıymetli danışmanımız Ömür Şahin'e sonsuz teşekkürlerimizi sunarız.

Kullanıcıların anlık ruh hallerini ve duygusal bağlamlarını anlayarak, müzik keşfini sıradan bir listeleme mantığından çıkarıp etkileşimli bir deneyime dönüştürme fikri, mevcut platformların kısıtlamalarından duyulan bir eksiklikten doğmuştur. Bu süreçte yapay zeka, doğal dil işleme ve mobil uygulama mimarileri üzerine yoğunlaşarak sadece bir bitirme projesi değil, gerçek dünyada kullanılabilir, ölçeklenebilir bir ürün vizyonu geliştirdik.

Eğitim hayatımız boyunca bize maddi ve manevi her türlü desteği sağlayan, zorlu dönemlerde yanımızda olan sevgili ailelerimize en derin şükranlarımızı sunarız. Ayrıca projede kullandığımız açık kaynak teknolojilere katkıda bulunan tüm geliştirici topluluklarına teşekkür ederiz.

\vspace{1cm}
\begin{flushright}
Berke BALCI\\
İlayda BARAN\\
Haziran 2026
\end{flushright}

% ── ÖZET ────────────────────────────────────────────────────
\newpage
\phantomsection
\addcontentsline{toc}{chapter}{ÖZET}
\chapter*{ÖZET}

\begin{center}
  {\large\bfseries MUSAİC: YAPAY ZEKA DESTEKLİ DUYGU DUYARLI MÜZİK ÖNERİ SİSTEMİ}\\[0.5cm]
  Berke Balcı, İlayda Baran\\[0.3cm]
  Erciyes Üniversitesi, Bilgisayar Mühendisliği\\[0.3cm]
  Danışman: Ömür Şahin
\end{center}

\vspace{0.5cm}
Günümüz dijital müzik platformları, kullanıcılara müzik önerisi sunarken ağırlıklı olarak geçmiş dinleme istatistiklerine dayalı işbirlikçi filtreleme yöntemlerini kullanmaktadır. Ancak bu geleneksel yaklaşım, kullanıcının o anki duygu durumunu ve spesifik bağlamını yakalamakta yetersiz kalmaktadır. Bu çalışmanın temel problemi, statik geçmiş verilere dayanan müzik öneri sistemlerinin anlık duygusal ihtiyaçlara yanıt verememesidir.

Musaic, kullanıcıların ruh hali metinlerini Google Gemini API ile anlamlandıran ve elde edilen verileri PyTorch tabanlı derin öğrenme modeli (Öklid mesafesi) ile 88.000 şarkılık veri setinde eşleştiren yenilikçi bir hibrit yapay zeka mimarisine dayanmaktadır. Arka uçta FastAPI ile yönetilen bu karmaşık algoritmalar, iOS platformuna geliştirilen kullanıcı dostu arayüz ve kaydırma (swipe) mekanizmasıyla doğrudan entegre edilerek müzik keşfini interaktif ve kişiselleştirilmiş bir mobil deneyime dönüştürmektedir. Bu çalışma, büyük dil modelleri ile derin öğrenme disiplinlerinin mobil uygulamalarda nasıl başarılı bir şekilde harmanlanabileceğini kanıtlamaktadır.

\vspace{0.5cm}
\textbf{Anahtar Kelimeler:} Doğal Dil İşleme (NLP), 9-Boyutlu Duygu Vektörü, MVVM Mimarisi, FastAPI, Firebase Firestore, IOS, Mobil Uygulama

% ── ABSTRACT ────────────────────────────────────────────────
\newpage
\phantomsection
\addcontentsline{toc}{chapter}{ABSTRACT}
\chapter*{ABSTRACT}

\begin{center}
  {\large\bfseries MUSAIC: AI-POWERED EMOTION-AWARE MUSIC RECOMMENDATION SYSTEM}\\[0.5cm]
  Berke Balcı, İlayda Baran\\[0.3cm]
  Erciyes University, Department of Computer Engineering\\[0.3cm]
  Supervisor: Ömür Şahin
\end{center}

\vspace{0.5cm}
Today's digital music platforms predominantly rely on collaborative filtering methods based on historical listening statistics to provide music recommendations. However, this traditional approach falls short in capturing the user's immediate emotional state and specific context. The fundamental problem addressed in this study is the inability of static, history-based music recommendation systems to respond to instantaneous emotional needs.

To solve this problem, an iOS-based application named "Musaic" was developed. It analyzes mood texts written by users in natural language using artificial intelligence and recommends songs that best fit their current emotional state. The original aspect of this work is the transformation of user texts into a 9-dimensional musical emotion vector via Google Gemini API, and matching this vector within a database of 114,000 songs using PyTorch models involving Euclidean and Cosine distances.

Within the scope of the research, the backend architecture was designed with Python-based FastAPI, migrating all heavy AI workloads to the server, while the mobile client side implemented a Protocol-Oriented MVVM architecture using Swift and SwiftUI. Following a swipe action on the mobile app, liked songs are transmitted to Firebase Firestore, and retrieved via API when the user navigates to the Favorites screen. Furthermore, through Apple MusicKit integration, high-resolution cover images and 30-second audio previews of the recommended songs are concurrently presented. Consequently, Musaic has transformed music discovery into an interactive smart assistant format responding to user's real-time emotions.

\vspace{0.5cm}
\textbf{Keywords:} Natural Language Processing (NLP), 9-Dimensional Mood Vector, MVVM Architecture, FastAPI, Firebase Firestore.

% ── İÇİNDEKİLER ────────────────────────────────────────────
\newpage
\tableofcontents

% ── TABLOLAR LİSTESİ ────────────────────────────────────────
\newpage
\phantomsection
\addcontentsline{toc}{chapter}{TABLOLAR LİSTESİ}
\listoftables

% ── ŞEKİLLER LİSTESİ ────────────────────────────────────────
\newpage
\phantomsection
\addcontentsline{toc}{chapter}{ŞEKİLLER LİSTESİ}
\listoffigures

% ── KISALTMALAR ─────────────────────────────────────────────
\newpage
\phantomsection
\addcontentsline{toc}{chapter}{KISALTMALAR}
\chapter*{KISALTMALAR}

\begin{tabular}{ll}
  API  & : Application Programming Interface (Uygulama Programlama Arayüzü) \\[3pt]
  DI   & : Dependency Injection (Bağımlılık Enjeksiyonu) \\[3pt]
  DTO  & : Data Transfer Object (Veri Transfer Nesnesi) \\[3pt]
  EMA  & : Exponential Moving Average (Üstel Hareketli Ortalama) \\[3pt]
  iOS  & : iPhone Operating System \\[3pt]
  JSON & : JavaScript Object Notation \\[3pt]
  LLM  & : Large Language Model (Büyük Dil Modeli) \\[3pt]
  MLP  & : Multilayer Perceptron (Çok Katmanlı Algılayıcı) \\[3pt]
  MVVM & : Model-View-ViewModel \\[3pt]
  NLP  & : Natural Language Processing (Doğal Dil İşleme) \\[3pt]
  REST & : Representational State Transfer \\[3pt]
  UI   & : User Interface (Kullanıcı Arayüzü) \\
\end{tabular}

% ════════════════════════════════════════════════════════════
%  ANA METİN — Arap rakamı numaralandırma başlar
% ════════════════════════════════════════════════════════════
\newpage
\pagenumbering{arabic}
\setcounter{page}{1}

% ────────────────────────────────────────────────────────────
%  GİRİŞ  (numarasız, zorunlu)
% ────────────────────────────────────────────────────────────
\chapter*{GİRİŞ}
\phantomsection
\addcontentsline{toc}{chapter}{GİRİŞ}
\markboth{GİRİŞ}{GİRİŞ}

Müzik, insanların duygu durumlarını yönetmek, ifade etmek ve ruh hallerini dengelemek için başvurdukları en evrensel araçtır. Dijitalleşme ile birlikte müzik tüketim alışkanlıkları tamamen değişmiş; milyonlarca şarkıya anında erişim imkânı, "ne dinleyeceğine karar verme" sürecini bir tür bilişsel yüke dönüştürmüştür. 

\vspace{0.5cm}
\noindent\textbf{Konunun Önemi ve Problem Tanımı}

Günümüzde Spotify, Apple Music gibi popüler müzik platformları, kullanıcılara müzik önerisi sunarken geçmiş dinleme alışkanlıklarına dayalı algoritmalar kullanmaktadır [1]. İnsanlar, yıllar boyunca oluşturdukları statik çalma listeleri veya sistemin "Bunu dinleyenler bunu da dinledi" mantığıyla ürettiği mekanik tavsiyeler üzerinden müzik keşfetmektedir. 

Ancak bu yaklaşımda büyük bir eksiklik vardır: İnsanın duygu durumu statik değildir. Kullanıcının o anki spesifik duygu durumuna (mood) uygun bir parça arayışı, geleneksel sistemler tarafından geçmiş verilerin gürültüsü sebebiyle yakalanamamaktadır. Musaic projesini geliştirme nedenimiz tam olarak bu boşluğu doldurmak, müzik platformlarını geçmişe dayalı birer arşivci olmaktan çıkarıp, anı yaşayan empatik birer asistan dönüştürmektir.

\vspace{0.5cm}
\noindent\textbf{Projenin Genel İşleyiş Süreci}

Musaic sisteminin genel işleyişi şu temel akış üzerine inşa edilmiştir:
1. Kullanıcı uygulamaya giriş yapar (Login) ve ana sayfa üzerinden sohbet ekranına yönelir.
2. Uygulama arayüzüne "Gece yolda gidiyormuş gibi hissediyorum" gibi doğal bir dilde serbest metin girilir.
3. iOS uygulaması, metni FastAPI tabanlı backend sunucusuna iletir.
4. Sunucudaki Google Gemini entegreli NLP motoru, metni analiz ederek 9 boyutlu matematiksel bir duygu vektörü üretir ve cihazla paylaşır.
5. Vektörel veriyle eşleşen en uygun şarkılar PyTorch modeli üzerinden hesaplanarak mobil cihaza iletilir.
6. Apple Music API üzerinden şarkıların yüksek çözünürlüklü kapak görselleri ve 30 saniyelik ses dosyaları mobil uygulamada zenginleştirilir.
7. Kullanıcı Tinder tarzı bir arayüzde (Swipe) şarkıları sağa kaydırarak (beğenerek) Firestore veritabanına ekler.

\vspace{0.5cm}
\noindent\textbf{Çalışmanın Amacı ve Ölçülebilir Hedefleri}

Bu çalışmanın **amacı**; müzik keşfini kullanıcının anlık ruh haliyle tam bir uyum içine sokarak, dijital müzik tüketiminde makine-insan etkileşimini doğal dil üzerinden kişiselleştirilmiş bir boyuta taşımaktır.

Proje kapsamında başarıyla tamamlanan **ölçülebilir hedefler** ise şunlardır:
\begin{itemize}
  \item Doğal dil metinlerini 9 boyutlu müzikal parametre vektörlerine dönüştürebilen bir NLP aracı katmanı geliştirmek.
  \item 88.000'den fazla kayıt barındıran veri seti üzerinde PyTorch destekli, milisaniyeler seviyesinde vektörel eşleştirme yapan bir Python tabanlı FastAPI servisi kurmak.
  \item Sunucu ile Swift arayüzü arasında 8 farklı ekrandan oluşan asenkron bir mobil istemci mimarisi inşa etmek.
  \item Swiped (kaydırılmış) verilerin Firebase Firestore üzerinde saklanması ve çekilmesini sağlayan bir altyapı oluşturmak.
\end{itemize}

\vspace{0.5cm}
\noindent\textbf{Tezin Organizasyonu}

Bu raporun ilerleyen bölümleri şu şekilde organize edilmiştir: 1. Bölümde, projenin arka planındaki yapay zeka altyapısı, veri seti ön işlemleri ve çift aşamalı tavsiye motoru teknik detaylarıyla incelenmiştir. 2. Bölümde, mobil uygulama mimarisi, MVVM deseni, klasör yapısı, API istek döngüleri ve kullanıcı ekranları detaylandırılmıştır. Son bölümde ise elde edilen sonuçlar değerlendirilmiş ve projenin limitasyonları ele alınmıştır.

% ────────────────────────────────────────────────────────────
%  1. BÖLÜM — YAPAY ZEKA VE BACKEND SİSTEMLERİ (Tez docx İçeriği)
% ────────────────────────────────────────────────────────────
\chapter{YAPAY ZEKA VE BACKEND SİSTEMLERİ}

Bu çalışmada geliştirilen müzik öneri motoru, modern veri bilimi literatüründe yüksek performanslı sistemlerin standardı olarak kabul edilen "İki Aşamalı Tavsiye Sistemi" (Two-Stage Recommender System) mimarisi üzerine inşa edilmiştir.

Sistem; kullanıcının doğal dil ile ifade ettiği soyut anlık ruh halini yakalamak, bu ruh haline matematiksel olarak en yakın şarkı adaylarını bulmak ve son olarak bu adayları derin öğrenme süzgecinden geçirerek en kaliteli sıralamayı üretmek üzere yapılandırılmıştır. Mimari; Büyük Dil Modelleri (LLM), Vektör Uzayı Modellemesi ve Derin Öğrenme katmanlarının melez (hibrit) bir formda çalıştırılması prensibine dayanır.

\section{Kullanılan Teknolojiler ve Geliştirme Ortamı}

Bu çalışmada, yapay zeka destekli müzik öneri motorunun inşası ve mobil uygulama entegrasyonu süreçlerinde yararlanılan temel yazılım araçları ile teknolojiler aşağıda sınıflandırılmıştır:

\textbf{Programlama Dili ve Ortam}
\begin{itemize}
    \item \textbf{Python 3.x:} Veri analizi, matematiksel modelleme ve arka uç geliştirme.
    \item \textbf{Geliştirme ve test süreci:} Visual Studio Code (VS Code), Jupyter Notebook.
\end{itemize}

\textbf{Doğal Dil İşleme ve Yapay Zeka}
\begin{itemize}
    \item \textbf{Google Gemini 2.5 Flash API:} Doğal dil girdilerinin anlamsal analizi ve JSON formatında 9 boyutlu soyut duygu durum vektörüne dönüştürülmesi.
    \item \textbf{PyTorch:} Aday şarkıların yapısal kalite skorlarını hesaplayan Çift Kanallı Çok Katmanlı Algılayıcı (MLP) derin öğrenme mimarisi.
\end{itemize}

\textbf{Veri Seti Yönetimi ve İşleme}
\begin{itemize}
    \item \textbf{NumPy \& Pandas:} Yaklaşık 88.000 şarkılık Spotify veri setinin temizlenmesi, yüksek boyutlu matris formatında işlenmesi ve bellekte yönetilmesi.
    \item \textbf{Scikit-Learn:} Veri setinin özellik ölçeklendirmesi (Min-Max Normalizasyonu) ve görselleştirme için boyut indirgeme işlemleri (PCA).
\end{itemize}

\textbf{API Geliştirme ve Dağıtım}
\begin{itemize}
    \item \textbf{FastAPI:} Yüksek eşzamanlılık (concurrency) destekli RESTful servislerin oluşturulması.
    \item \textbf{Uvicorn:} Asenkron ASGI sunucu ortamı çalıştırıcısı.
    \item \textbf{Ngrok / Localhost.run:} Geliştirme ve test aşamasında iOS istemcisi ile yerel sunucu (localhost) arasında güvenli iletişim tünellerinin kurulması.
\end{itemize}

\textbf{Diğer}
\begin{itemize}
    \item \textbf{Git / GitHub:} Sürüm kontrolü ve işbirliği.
\end{itemize}

\section{Veri Seti Analizi ve Semantik Vektör Haritalama}

Geleneksel tavsiye sistemleri, genellikle kullanıcıların geçmiş dinleme kayıtlarına veya benzer kullanıcıların davranışlarına dayanır. Ancak bu durum, sistemi ilk kez kullanan kişiler için "Soğuk Başlangıç" (Cold-Start) problemi yaratır. Bu problemi çözmek için sistemin giriş katmanı, kullanıcıların anlık ruh hallerini serbest metin olarak ifade ettikleri Doğal Dil İşleme (NLP) modülü olarak tasarlanmıştır.

Klasik anahtar kelime eşleştirme yöntemleri soyut ifadelerde yetersiz kaldığından, sisteme semantik (anlamsal) yönlendirici olarak Google Gemini 2.5 Flash modeli entegre edilmiştir.

\subsection{Hedef Vektör Üretimi ve Veri Dağılımı}
Geliştirilen komut mühendisliği (prompt engineering) yapısı sayesinde kullanıcının girdiği ham metin (örn. "yağmurlu havada dinlenecek sakin şarkılar"), Spotify altyapısından elde edilen 88.871 geçerli şarkının akustik DNA'sını temsil eden 9 boyutlu sürekli bir hedef vektöre dönüştürülür (Danceability, Energy, Valence, Tempo, Acousticness, Instrumentalness, Speechiness, Loudness, Liveness). Modelin verileri matematiksel olarak daha adil işleyebilmesi için veri setindeki akustik özelliklerin dağılımları ve birbirleriyle olan ilişkileri analiz edilmiştir:

% RESİM EKLEME NOKTASI: Dağılım Grafikleri için
\begin{figure}[H]
    \centering
     \includegraphics[width=\linewidth]{dagilim_grafigi.png}
    \vspace{6cm} % Resim eklendiğinde bu vspace komutunu siliniz.
    \caption{Müzik veri setinde yer alan 9 temel akustik özniteliğin dağılım grafikleri.}
    \label{fig:dagilim}
\end{figure}

% RESİM EKLEME NOKTASI: Korelasyon Isı Haritası için
\begin{figure}[H]
    \centering
     \includegraphics[width=0.85\linewidth]{korelasyon_isi.png}
    \vspace{6cm} % Resim eklendiğinde bu vspace komutunu siliniz.
    \caption{Akustik öznitelikler arasındaki istatistiksel ilişkileri gösteren korelasyon ısı haritası.}
    \label{fig:isi_haritasi}
\end{figure}

\section{Aday Üretimi (Retrieval) ve Mesafe Algoritmaları}

Veritabanındaki on binlerce şarkının tek tek ağır yapay zeka modellerinden geçirilmesi sistemin yavaşlamasına neden olur. Bu nedenle, LLM tarafından üretilen 9 boyutlu hedef vektör ile veritabanındaki şarkılar arasında hızlı bir ön eleme (Retrieval) yapılması gerekmiştir.

\subsection{Kosinüs Benzerliği Yanılsaması ve Öklid Çözümü}
Geliştirme sürecinin ilk aşamasında, tavsiye sistemlerinde standart olarak kabul edilen Kosinüs Benzerliği kullanılmıştır. Ancak testler sırasında önemli bir matematiksel sorun fark edilmiştir: Kosinüs benzerliği, vektörlerin büyüklüklerine değil, aralarındaki açıya odaklanmaktadır.

Örneğin, kullanıcının çok düşük enerji ve düşük valans (hüzün) talep ettiği senaryolarda, sadece özellik oranları birbirine benziyor diye sistemin yüksek enerjili ve gürültülü pop şarkılarını da önerdiği tespit edilmiştir. Bu "Açısal Oran Yanılsamasını" ortadan kaldırmak için algoritma güncellenerek, uzaydaki mutlak mesafeyi ölçen **Öklid Mesafesi (Euclidean Distance)** formülüne geçilmiştir. Bu sayede kullanıcının istediği duygunun "şiddeti" tam isabetle yakalanmıştır.

\subsection{Dinamik Min-Max Ölçekleme}
Öklid mesafesine geçildikten sonra, aşırı uç duygu durumlarında (örn. tamamen akustik bir talepte) hedefe yakın hiçbir şarkı kalmadığında mesafelerin çok açılmasıyla yaşanan "Sıfır Benzerlik / Boş Liste" problemi ortaya çıkmıştır. Bu durumu çözmek için uzaklıklara Dinamik Min-Max Ölçekleme uygulanmıştır:

$$S_{\text{Euclidean}} = 1.0 - \left( \frac{d - d_{\min}}{d_{\max} - d_{\min}} \right)$$

Bu bağıl değerlendirme algoritması sayesinde sistem, mutlak mesafeler ne kadar uzak olursa olsun, uzaydaki en yakın komşuları bularak o anki veritabanının en iyi adayına bağıl olarak maksimum uyum oranını (\%100) atamayı başarmıştır.

% RESİM EKLEME NOKTASI: PCA ile Boyut İndirgeme Grafiği için
\begin{figure}[H]
    \centering
     \includegraphics[width=\linewidth]{Picture1.png}
    \vspace{6cm} % Resim eklendiğinde bu vspace komutunu siliniz.
    \caption{PCA ile Boyut İndirgeme ve Öklid Mesafe Simülasyonu.}
    \label{fig:pca_grafik}
\end{figure}

\section{Çift Kanallı Yapay Sinir Ağı ile Nihai Sıralama (Ranking)}

Öklid mesafe filtresinden geçmeyi başaran ve kullanıcı moduna en uygun 100 şarkılık rafine aday havuzu, sistemin en zeki katmanı olan PyTorch tabanlı Çok Katmanlı Algılayıcı (MLP) modeline beslenir. Bu derin öğrenme modelinin amacı, şarkıları sadece yüzeydeki özelliklerine göre değil, müzikal yapı kalitesine göre puanlamaktır.

Model, çift kanallı (dual-input) bir yapı ile çalışır:
\begin{itemize}
    \item \textbf{Gömme (Embedding) Kanalı:} Şarkıların kimlik numaraları (song\_id), 32 boyutlu sürekli bir gizli uzaya izdüşürülerek birbirleriyle olan yapısal ilişkileri (hangi şarkı hangisiyle benzer) öğrenilir.
    \item \textbf{Akustik Kanal:} 9 boyutlu ses özellikleri, doğrusal olmayan karmaşık örüntülerin çözülmesi için ReLU aktivasyonlu katmanlardan geçirilerek 32 boyutlu bir matrise dönüştürülür.
\end{itemize}

Çıkış katmanındaki Sigmoid aktivasyon fonksiyonu bu birleştirilmiş veriyi işleyerek şarkının yapısal kalitesini temsil eden bir kalite skoru üretir.

\textbf{Performans Optimizasyonu (Caching):} Sistem performansını yüksek tutmak ve yanıt süresini (latency) düşürmek amacıyla bu ağır derin öğrenme skorları, API sunucusu ayağa kalktığında bir kez hesaplanıp önbelleğe (caching) alınır. Canlı algoritmada bu baz skorlar, anlık Öklid benzerlik oranlarıyla harmanlanarak nihai tavsiye puanını oluşturur:

$$Score_{\text{Nihai}} = (0.15 \times Score_{\text{PyTorch}}) + (0.85 \times S_{\text{Euclidean}})$$

Bu formül, kullanıcının o anki anlık ruh haline sıkı sıkıya uyulmasını sağlarken (\%85), müzikal kalitesi çok düşük yapımların listeye girmesini (\%15) engellemektedir.

\section{Durumsuz (Stateless) Etkileşim ve Gerçek Zamanlı Öğrenme}

Kullanıcının zaman içerisindeki zevk değişimlerine adapte olabilmek için oyunlaştırılmış bir "Kaydırma (Swipe)" mekanizması tasarlanmıştır. Sunucu üzerindeki veritabanı okuma/yazma (I/O) yükünü minimuma indirmek adına sistem tamamen Durumsuz (Stateless) mimaride kurgulanmıştır. Kaydırma etkileşimleri sunucu veritabanına anlık yazılmaz; state (durum) yönetimi doğrudan iOS istemcisinin yerel belleğinde yürütülür.

Kullanıcının sağa kaydırarak beğendiği şarkıların akustik özellikleri, istatistiksel bir yöntem olan Üstel Hareketli Ortalama (EMA) mantığıyla kullanıcının oturum vektörüne etki eder:

$$V_{\text{yeni}} = V_{\text{mevcut}} + \alpha \cdot (\mu_{\text{beğenilen}} - V_{\text{mevcut}})$$

Bu matematiksel yaklaşım; kullanıcının zevklerini baştan sıfırlamak yerine, yeni beğendiği şarkılara ($\mu$) doğru adım adım yumuşak bir şekilde kaydırmaktadır. Güncellenen bu vektör bir sonraki istekte sunucuya yük (payload) olarak taşınarak, ağır model eğitimlerine (retraining) gerek kalmadan gerçek zamanlı (real-time) adaptasyon sağlar.

\section{Son İşleme (Post-Processing) ve Çeşitlilik Mekanizmaları}

Kullanıcı arayüzüne (UI) gönderilecek nihai 15 şarkılık listenin veri kalitesini artırmak ve kullanıcı deneyimini taze tutmak için üç aşamalı bir güvenlik ve çeşitlilik protokolü uygulanır:

\begin{itemize}
    \item \textbf{Kopya Veri Filtresi (Deduplication):} Spotify veri setinde sıkça rastlanan, aynı şarkıya ait mükerrer kayıtlar (Single ve Albüm varyasyonları), sıralama aşamasında şarkı ismi ve sanatçı bazında temizlenir. Böylece kullanıcının aynı şarkıyı iki kez görmesi engellenir.
    \item \textbf{Popülerlik Dengesi (Popularity Bias Mitigation):} Tavsiye algoritmalarının her zaman en çok dinlenen hit şarkılara yönelme eğilimini kırmak için, popülerlik çarpanının genel skora etkisi maksimum \%10 seviyesine kalibre edilmiştir.
    \item \textbf{Ağırlıklı Keşif Örneklemesi (Weighted Random Sampling):} En yüksek puan alan şarkıları her zaman statik bir liste olarak sunmak (sömürü/exploitation), sistemin yeni parçalar keşfettirme (keşif/exploration) özelliğini yok etmektedir. Bunun çözümü olarak sistem, en iyi 100 şarkılık geniş bir aday havuzu oluşturur. Bu havuzdan, şarkıların final skorları birer çekiliş ihtimali (weights) olarak kullanılarak rastgele 15 şarkı örneklenir. Bu sayede uygulama, müzikal isabet oranından ödün vermeden her oturumda kullanıcıya yepyeni ve sürpriz parçalar sunmayı başarmaktadır.
\end{itemize}

\section{API Geliştirme ve Sunucu Mimarisi Entegrasyonu}

Bu çalışmada tasarlanan hibrit müzik öneri motorunun (doğal dil işleme, vektör uzayı hesaplamaları ve derin öğrenme katmanları) son kullanıcıya ulaştırılabilmesi için istemci-sunucu (client-server) mimarisine dayalı bir entegrasyon stratejisi benimsenmiştir.

Yaklaşık 88.000 satırlık devasa veri matrisinin (\texttt{X\_scaled.npy}), hesaplama maliyeti yüksek PyTorch modelinin ve API anahtarlarının doğrudan uç cihazda (edge device / iOS) konumlandırılması, mobil cihazlarda bellek darboğazlarına (memory bottlenecks), performans kayıplarına ve güvenlik zafiyetlerine yol açacaktır. Bu bağlamda, mobil cihaz yalnızca bir sunum ve etkileşim katmanı (presentation layer) olarak bırakılmış; tüm bilişsel hesaplama yükü Python tabanlı bir arka uç (backend) sunucusuna devredilmiştir.

İstemci ile sunucu arasındaki asenkron veri akışını koordine etmek amacıyla RESTful mimari standartlarına uygun olarak \textbf{FastAPI} çerçevesi (framework) kullanılmıştır. ASGI tabanlı \textbf{Uvicorn} sunucusu üzerinde koşan bu yapı, donanım kaynaklarını optimize etmek adına durumsuz (stateless) bir iletişim protokolü ile aşağıdaki uç noktalar üzerinden hizmet vermektedir:

\begin{itemize}
    \item \textbf{POST /api/chat:} Doğal dil girdisini semantik olarak ayrıştıran ve LLM (Google Gemini) çağrılarını yöneten uç noktadır. Gelen soyut metin dizgisi analiz edilerek, 9 boyutlu sürekli bir duygu durum vektörüne haritalanır ve JSON yükü (payload) formunda istemciye iletilir.
    \item \textbf{POST /api/recommend:} Aday üretimi (retrieval) ve nihai sıralama (ranking) algoritmalarının tetiklendiği merkezi modüldür. İstemciden gelen hedef vektörü kullanılarak tüm veri seti üzerinde Öklid mesafesi hesaplanır, önbellekteki (cache) Çok Katmanlı Algılayıcı (MLP) skorlarıyla harmanlanır ve dışlama filtreleri (deduplication) uygulandıktan sonra en uygun 15 şarkılık liste (önizleme bağlantıları ile birlikte) döndürülür.
    \item \textbf{POST /api/swipe:} Kullanıcının arayüzdeki ikili etkileşimlerini (sağa/sola kaydırma) işleyen modüldür. Kullanıcı zevklerini öğrenmek adına, Üstel Hareketli Ortalama (EMA) yöntemi ile kullanıcının profil vektörünü kümülatif olarak güncelleyerek gerçek zamanlı adaptasyon sağlar.
\end{itemize}

% RESİM EKLEME NOKTASI: Swagger UI /api/chat ve /api/recommend
\begin{figure}[H]
    \centering
    % \includegraphics[width=\linewidth]{swagger_api_endpoints.png}
    \vspace{6cm} % Resim eklendiğinde bu vspace komutunu siliniz.
    \caption{Swagger UI (OpenAPI) arayüzünde Musaic RESTful API uç noktalarının yapısal görünümü.}
    \label{fig:swagger_endpoints}
\end{figure}

\subsection{Asenkron G/Ç Yönetimi ve Zaman Aşımı İyileştirmeleri}
Ağ tabanlı gecikmelerin (network latency) ve LLM çıkarım (inference) sürelerinin istemci tarafında uygulamayı kilitlemesini (thread-blocking) önlemek elzemdir. Bu doğrultuda, tüm API uç noktaları eşzamanlı girdi/çıktı (I/O) işlemlerine izin verecek şekilde asenkron (async/await) bloklar halinde kodlanmıştır.

Geliştirme safhasında Gemini modelinin yanıt süresinin uzamasıyla karşılaşılan ağ zaman aşımı (Timeout: -1001) hataları; sunucu tarafındaki asenkron kuyruk yönetiminin optimize edilmesi ve eş zamanlı olarak iOS istemcisindeki zaman aşımı eşik değerlerinin (timeout intervals) yeniden kalibre edilmesiyle kararlı (stable) hale getirilmiştir.

\subsection{Geliştirme Ortamı, Tünelleme ve Test Süreçleri}
Sistemin geliştirme ve doğrulama (validation) süreçlerinde, sunucu ortamı yerel donanım (localhost) üzerinde ayağa kaldırılmıştır. Mobil istemcinin (iOS cihazı) bu yerel sunucuya dış ağlar üzerinden güvenli ve kesintisiz erişim sağlayabilmesi için SSH tünelleme çözümlerinden (Örn. \textbf{Ngrok} veya \textbf{Localhost.run}) faydalanılmıştır. Bu yaklaşım, sistemin canlı üretim ortamı (production environment) davranışlarını birebir simüle etmesini sağlarken, aynı zamanda hızlı prototipleme ve gerçek zamanlı hata ayıklama (debugging) imkânı sunmuştur.

% RESİM EKLEME NOKTASI: Swagger UI JSON Payload
\begin{figure}[H]
    \centering
    % \includegraphics[width=\linewidth]{swagger_json_payload.png}
    \vspace{6cm} % Resim eklendiğinde bu vspace komutunu siliniz.
    \caption{Swagger arayüzü üzerinden API tahmin sonuçlarının (JSON payload) test edilmesi ve doğrulanması.}
    \label{fig:swagger_payload}
\end{figure}

API uç noktalarının yapısal analizi, sözleşme (contract) testleri ve veri gövdesi (payload) doğrulamaları için FastAPI'nin entegre \textbf{Swagger UI (OpenAPI)} belgelendirme arayüzü aktif olarak kullanılmıştır. Bu durumsuz (stateless) ve asenkron mimari; sunucu üzerinde oturum (session) tutma zorunluluğunu ve veritabanı darboğazlarını ortadan kaldırarak projenin gereksinimlerine uygun, düşük gecikmeli (low-latency) ve ölçeklenebilir bir mühendislik çözümü ortaya koymuştur.


% ────────────────────────────────────────────────────────────
%  2. BÖLÜM — MOBİL UYGULAMA GELİŞTİRME
% ────────────────────────────────────────────────────────────
\chapter{MOBİL UYGULAMA GELİŞTİRME}

Projenin kullanıcı ile etkileşime giren istemci (client) katmanı, Apple ekosisteminin en güncel teknolojileri kullanılarak iOS platformu için geliştirilmiştir. Bu bölümde, uygulamanın yazılım mimarisi, tercih edilen teknolojilerin gerekçeleri, sistem içi veri akışları ve kullanıcı arayüzü tasarımları detaylandırılmaktadır.

\section{Mobil Uygulama Mimarisi ve Teknoloji Seçimi}

Musaic iOS uygulaması, sürdürülebilir, test edilebilir ve ölçeklenebilir bir yapı sunmak amacıyla \textbf{MVVM (Model-View-ViewModel)} mimari deseni kullanılarak yapılmıştır. Arayüz geliştirmeleri için modern, deklaratif bir framework olan SwiftUI tercih edilmiştir.
\subsection{Uygulama Genel Akışı}

Kullanıcının uygulamaya girişiyle başlayan ve müzik keşif süreciyle devam eden genel sistem akışı, kesintisiz bir kullanıcı deneyimi sunacak şekilde tasarlanmıştır. Şekil \ref{fig:app_flow}'da uygulamanın temel modülleri arasındaki etkileşim ve yönlendirme akışı gösterilmektedir. 

% RESİM EKLEME NOKTASI: UYGULAMA GENEL AKIŞI
\begin{figure}[H]
    \centering
    \vspace{-0.5cm}
\includegraphics[width=0.95\linewidth, height=10cm]{User Login Flow for Chat-2026-06-20-105344.png}
     % Yüksekliği 15 cm yapar
    \vspace{-0.3cm}
    \caption{Uygulama Genel Akış Diyagramı (Flowchart).}
    \label{fig:app_flow}
\end{figure}

\subsection{Protokol Odaklı Bağımlılık Enjeksiyonu (Protocol-Oriented Dependency Injection)}

Uygulamanın servis katmanında dışa bağımlılıkları yönetmek amacıyla Protokol Odaklı Bağımlılık Enjeksiyonu kullanılmıştır. Üçüncü parti bir kütüphane yerine tamamen Swift'in yerleşik özellikleri ile tasarlanan manuel bir DI (Dependency Injection) konteyneri oluşturulmuştur. Bu yaklaşımda, servis sınıfları doğrudan kullanılmak yerine (örneğin \texttt{APIRecommendationService}), öncelikle bir protokol (örneğin \texttt{RecommendationServiceProtocol}) üzerinden soyutlanmıştır. \texttt{DIContainer} sınıfı, bu protokolleri uygulayan somut sınıfların (concrete classes) yaşam döngülerini yönetir. Geliştirme sürecinde gerçek API servisleri yerine sahte (mock) servislerin sisteme entegre edilebilmesi bu sayede kolaylaşmıştır.

% RESİM EKLEME NOKTASI: MİMARİ DİYAGRAMI
\begin{figure}[H]
    \centering
    \vspace{-0.5cm}
     \includegraphics[width=0.9\linewidth]{SistemDIMimarisi.png}
    \vspace{-0.3cm}
    \caption{Mobil Uygulama Katmanlı Mimarisi ve Servis İletişimi Diyagramı.}
    \label{fig:architecture_diagram}
\end{figure}

Bu mimarinin merkezinde, mobil istemcinin ağır bilişsel yüklerden arındırılarak "ince istemci" (thin client) prensibiyle çalışması yatmaktadır. Tüm karmaşık model hesaplamalarının sunucu tarafında işlenmesi sayesinde uygulamanın hem derleme boyutu düşürülmüş hem de işlemci/batarya optimizasyonu sağlanmıştır:

\begin{itemize}
    \item \textbf{Doğal Dil İşleme için Gemini API Kullanımı:} Kullanıcının yazdığı serbest metinlerden 9 boyutlu duygu vektörleri çıkarmak karmaşık bir anlamsal eşleştirme gerektirir. Bu işlemi cihaz üzerinde (on-device) çalışan yerel bir modelle yapmak, uygulama boyutunu gigabaytlarca artıracak ve cihaz bataryasını hızla tüketecektir. Bu nedenle metin analizi, yüksek hızlı ve hafif bir LLM olan \textit{Google Gemini API} kullanılarak FastAPI sunucusu tarafında gerçekleştirilmiş, mobil cihaz sadece giriş metnini iletip JSON formatındaki vektörü alan bir aracı olarak konumlandırılmıştır.
    \item \textbf{Öneri Motoru için FastAPI Aracılığıyla PyTorch İstekleri:} Yaklaşık 88.000 şarkıdan oluşan bir vektör uzayında hesaplama yapmak ve derin öğrenme modelini çalıştırmak iOS cihazlar için bellek darboğazlarına (memory bottlenecks) yol açma riskine sahiptir. Bu sebeple geliştirilen PyTorch modeline asenkron işlemler sunan \textbf{FastAPI} aracılığıyla istek (request) atılmaktadır. Mobil uygulama, HTTP POST istekleri ile bu servise bağlanarak kısa bir sürede filtrelenmiş şarkı sonuçlarını elde etmektedir.
\end{itemize}

Aşağıda, mimari tasarımı ve proje modüllerini açıkça yansıtan temel klasör yapısı (Folder Structure) sunulmuştur.

% RESİM EKLEME NOKTASI: KLASÖR YAPISI
\begin{figure}[H]
    \centering
    \vspace{-0.5cm}
    \includegraphics[width=0.6\linewidth]{klasoryapisi.png}
    \vspace{-0.3cm}
    \caption{Musaic iOS Projesinin MVVM mimarisini gösteren Klasör Yapısı.}
    \label{fig:folder_structure}
\end{figure}

\section{Sistem İçi İletişim ve Veri Akışı (Sequence)}

Mobil istemci ile FastAPI sunucusu arasındaki iletişim, asenkron bir JSON veri akışına dayanır. Uygulamanın çalışması, iOS cihazı ile harici servisler arasında sırasıyla şu adımları izler:

\begin{enumerate}
    \item \textbf{iOS Uygulaması} $\rightarrow$ \textbf{FastAPI} (\texttt{/api/chat} isteği - Kullanıcı metni gönderilir)
    \item \textbf{FastAPI} $\rightarrow$ \textbf{Gemini API} (Metin analiz edilir)
    \item \textbf{Gemini API} $\rightarrow$ \textbf{FastAPI} (9D Vektör üretilir)
    \item \textbf{FastAPI} $\rightarrow$ \textbf{iOS Uygulaması} (Vektör iOS'a döner)
    \item \textbf{iOS Uygulaması} $\rightarrow$ \textbf{FastAPI} (\texttt{/api/recommend} isteği - Vektör gönderilir)
    \item \textbf{FastAPI} $\rightarrow$ \textbf{PyTorch Modeli} (Vektör eşleştirmesi / Cosine Similarity yapılır)
    \item \textbf{PyTorch Modeli} $\rightarrow$ \textbf{FastAPI} (Şarkı ID'leri bulunur)
    \item \textbf{FastAPI} $\rightarrow$ \textbf{iOS Uygulaması} (Şarkı listesi döner)
    \item \textbf{iOS Uygulaması} $\rightarrow$ \textbf{Apple MusicKit} (Şarkılar zenginleştirilir - Kapak fotoğrafı, 30sn ses)
\end{enumerate}

Sistemde kullanici duygu durumunu girdikten sonra soyle bir akış olur:



% RESİM EKLEME NOKTASI: REQUEST/RESPONSE MODEL
\begin{figure}[H]
    \centering
    \vspace{-0.5cm}
     \includegraphics[width=0.8\linewidth]{moodinput_sequence.png}
    \vspace{-0.3cm}
    \caption{Sistemin kullaniciya sarki onerme mekanizmasi sequence diyagrami.}
    \label{fig:req_res_model}
\end{figure}

\section{Kullanıcı Akışı ve Arayüz Tasarımı}

Musaic uygulamasında özgün ekran tasarlanmış olup, kullanıcı arayüzü modern ve minimalist bir estetik sunan "Glassmorphism" tasarım dili etrafında şekillendirilmiştir. Kullanıcı deneyimi kesintisiz bir yapı üzerine kurulmuştur.

\textbf{1. Kayıt Olma (Sign Up) Ekranı:} Sisteme ilk kez giriş yapacak kullanıcılar için tasarlanan Kayıt Olma ekranı, kullanıcı deneyimini zorlaştırmayan sade bir form yapısına sahiptir. Güvenli iletişim için Firebase Authentication altyapısı kullanılmış olup, kullanıcıdan temel e-posta ve şifre bilgileri alınarak sisteme kaydı gerçekleştirilir.

% RESİM EKLEME NOKTASI: SIGN UP EKRANI
\begin{figure}[H]
    \centering
    \vspace{-0.5cm}
     \includegraphics[width=0.4\linewidth]{kayitOlma_page.png}
    \vspace{-0.3cm}
    \caption{Kayıt Olma (Sign Up) ekranının tasarımı.}
    \label{fig:signup_screen}
\end{figure}

\textbf{2. Login ve Kimlik Doğrulama Ekranı:} Kayıtlı kullanıcıların oturum açmasını sağlayan Login ekranı, uygulamanın açılışında karşılanan ilk güvenlik katmanıdır. Firebase Authentication üzerinden başarıyla doğrulanan kullanıcı oturumu, uygulama içi API isteklerinde kullanılmak üzere güvenli bir ID token sağlar. Kayitli kullanicilarinin karsilarina Kayit olma sayfasi cikmaz

% RESİM EKLEME NOKTASI: LOGIN EKRANI
\begin{figure}[H]
    \centering
    \vspace{-0.5cm}
    \includegraphics[width=0.4\linewidth]{login_ekrani.png}
    \vspace{-0.3cm}
    \caption{Kullanıcı Giriş (Login) ekranının tasarımı.}
    \label{fig:login_screen}
\end{figure}

\textbf{3. Ana Sayfa (Landing Home) ve Alt Navigasyon Çubuğu:} Giriş işlemi başarıyla tamamlandıktan sonra kullanıcı, \texttt{LandingHomeView} ekranına yönlendirilir. Bu ekran, kullanıcının dikkatini çeken dinamik bir karşılama (Hero Banner) bölümüne ve uygulamanın temel işlevlerine hızlı erişim sağlayan yönlendirme butonlarına ev sahipliği yapar. Uygulama içi tüm ana gezinme ise ekranın alt kısmında bulunan özel tasarım \textit{Bottom Navigation Bar} ile yönetilir. Bu sayede kullanıcılar sekmeler arası kesintisiz ve hızlı bir geçiş yapabilirler.

% RESİM EKLEME NOKTASI: LANDING HOME EKRANI
\begin{figure}[H]
    \centering
    \vspace{-0.5cm}
     \includegraphics[width=0.4\linewidth]{landing_page.png}
    \vspace{-0.3cm}
    \caption{Ana Sayfa (Landing Home) ve Alt Navigasyon Çubuğu tasarımı.}
    \label{fig:landing_home_screen}
\end{figure}

\textbf{4. Chat (Duygu Analizi) Ana Ekranı:} Sistemin kalbi olan bu arayüzde kullanıcı, anlık duygu durumunu serbest metin olarak ifade eder. Mesaj iletildiğinde sistem \texttt{/api/chat} endpointine istek atarak arka planda vektör dönüşüm sürecini başlatır. Olusan Vektor ile /api/recommend endpointine tekrar istek atilir ve backend tarafinda egittimiz yapay zeka modeli tarafindan kullanicinin durumuna gore 5 adet sarki doner.

% RESİM EKLEME NOKTASI: CHAT EKRANI
\begin{figure}[H]
    \centering
    \vspace{-0.5cm}
     \includegraphics[width=0.4\linewidth]{chat_page.png}
    \vspace{-0.3cm}
    \caption{Kullanıcının metin girdiği ve Gemini destekli bot ile etkileşime girdiği Chat Ana Sayfası.}
    \label{fig:chat_screen}
\end{figure}

\begin{figure}[H]
    \centering
    \vspace{-0.5cm}
     \includegraphics[width=0.4\linewidth]{chat2ailoading_page.png}
    \vspace{-0.3cm}
    \caption{Kullanıcının metin girdiği ve Gemini destekli bottan yanit bekledigi sayfanin goruntusu.}
    \label{fig:chat_screen}
\end{figure}

\begin{figure}[H]
    \centering
    \vspace{-0.5cm}
     \includegraphics[width=0.4\linewidth]{chat_ekrani.png}
    \vspace{-0.3cm}
    \caption{Kullanıcının girdigi metne gore backend'den donen sarkilarin oldugu ekran goruntusu.}
    \label{fig:chat_screen}
\end{figure}

\textbf{5. Swipe (Kaydırma) Arayüzü:} API üzerinden dönen şarkı tavsiyeleri, kullanıcının şarkıları Tinder benzeri oyunlaştırılmış bir arayüzle sağa (Beğen) veya sola (Reddet) kaydırabileceği şekilde sunulur. 

% RESİM EKLEME NOKTASI: SWIPE EKRANI
\begin{figure}[H]
    \centering
    \vspace{-0.5cm}
     \includegraphics[width=0.4\linewidth]{discovery_page.png}
    \vspace{-0.3cm}
    \caption{Önerilen şarkıların interaktif biçimde değerlendirildiği Swipe ekranı tasarımı.}
    \label{fig:swipe_screen}
\end{figure}

\begin{figure}[H]
    \centering
    \vspace{-0.5cm}
     \includegraphics[width=0.4\linewidth]{Swipe_page.png}
    \vspace{-0.3cm}
    \caption{Rastgele Sarki bulma Sayfasi.}
    \label{fig:swipe_screen}
\end{figure}

\section{Favoriler Yönetimi ve Veri Kalıcılığı}

Sistemde Swipe (kaydırma) işlemi bittikten sonra, beğenilen şarkılar ve bu şarkıların parametreleri (iOS tarafındaki \texttt{Song} sınıfının modeline göre) paket halinde Firebase Firestore veritabanına gönderilmektedir. 

Kullanıcı "Favorites" (Favoriler) ekranına geçtiğinde, cihaz içindeki sınırlı belleği zorlamamak adına bu veriler API aracılığıyla doğrudan Firestore'dan asenkron olarak çekilip listelenmektedir.

% RESİM EKLEME NOKTASI: FAVORITES EKRANI
\begin{figure}[H]
    \centering
    \vspace{-0.5cm}
     \includegraphics[width=0.4\linewidth]{favorites_page.png}
    \vspace{-0.3cm}
    \caption{Firestore üzerinden çekilen ve kullanıcının geçmiş beğenilerini listeleyen Favoriler Ekranı.}
    \label{fig:favorites_screen}
\end{figure}

\section{Apple MusicKit ile Medya Zenginleştirme}

FastAPI üzerinden gelen öneri sonuçları sadece veri setindeki temel metadata bilgilerini (şarkı adı, sanatçı adı, tür ve eşleşme skoru gibi) içerir. Mobil uygulama, Apple Music API (MusicKit) hizmetini entegre ederek listedeki şarkıların yüksek çözünürlüklü kapak görsellerini ve 30 saniyelik ses önizlemelerini cihaz üzerinden aratıp çeker. Bu özellik, uygulamanın müzikal sunum kalitesini en üst seviyeye taşımaktadır.

% ────────────────────────────────────────────────────────────
%  3. BÖLÜM — TARTIŞMA, SONUÇ VE ÖNERİLER
% ────────────────────────────────────────────────────────────
\chapter{TARTIŞMA, SONUÇ VE ÖNERİLER}

Bu bölümde projenin geldiği son nokta değerlendirilmiş, mevcut limitasyonlar şeffafça belirtilmiş ve gelecekte sisteme entegre edilmesi planlanan özellikler tartışılmıştır.

\section{Genel Değerlendirme ve Sonuçlar}

Musaic projesi; doğal dil işleme yeteneklerini ve derin öğrenme tabanlı müzik tavsiye algoritmalarını başarılı bir iOS uygulamasında birleştirmeyi başarmıştır. Kullanıcının ruh halini algılayıp 9 boyutlu bir müzikal vektör çıkarabilen altyapı, geleneksel statik dinleme algoritmalarının yetersizliklerine çözüm sunmuştur. Geliştirilen "İnce İstemci" mobil mimarisi, işlemlerin tüm ağırlığını FastAPI sunucusunda tutarak batarya ve performans dostu bir sistem yaratılmasını sağlamıştır.

\section{Projedeki Limitasyonlar}

Proje geliştirilirken donanımsal ve ticari ekosistem kuralları gereği karşılaşılan bazı limitasyonlar şu şekildedir:
\begin{itemize}
    \item \textbf{Medya Çalma Kısıtlamaları:} Müzik endüstrisindeki telif hakları nedeniyle Apple MusicKit API, aktif Apple Music aboneliği olmayan kullanıcılara sadece 30 saniyelik ses önizlemeleri sağlamaktadır. Bu durum uygulamanın tam bir medya oynatıcısı (streaming app) olmasını engellemiştir.
    \item \textbf{Model Çıkarım (Inference) Süresi:} Gemini API'sine atılan sorgular ve ardından gelen PyTorch vektör hesaplamaları ağ trafiği gecikmesine (network latency) tabidir. İşlemler sunucuda optimize edilmiş olsa da, yavaş bağlantılarda kullanıcının sonucu bekleme süresi artabilmektedir.
\end{itemize}

\section{Gelecek Çalışmalar ve Öneriler}

Projenin bir sonraki aşamalarında geliştirilmesi öngörülen başlıklar aşağıdadır:
\begin{itemize}
    \item Apple Watch (watchOS) entegrasyonu ile cihazdan anlık kalp atış hızı (nabız) gibi biyometrik verilerin alınması ve müzik duygu analizi modeline dahil edilmesi.
    \item Kullanıcıların kendi oluşturdukları duygu bazlı "Mood Playlist"lerini sosyal bir ağ mantığı ile diğer kullanıcılarla paylaşabilmesi (Mood Haritası).
    \item Çoklu dil desteğinin (Localization) genişletilmesi ve sadece Türkçe/İngilizce değil, diğer dillerdeki duygusal ifadelerin de modele daha hassas (fine-tuned) tanıtılması.
\end{itemize}

% ────────────────────────────────────────────────────────────
%  KAYNAKLAR 
% ────────────────────────────────────────────────────────────
\newpage
\phantomsection
\addcontentsline{toc}{chapter}{KAYNAKLAR}
\renewcommand{\bibname}{KAYNAKLAR}
\begin{thebibliography}{9}

\bibitem{ricci2011}
Ricci, F., Rokach, L. ve Shapira, B., 2011.
\textit{Introduction to Recommender Systems Handbook}.
Springer, Boston.

\bibitem{spotify2016}
Jacobson, K. et al., 2016.
Music Personalization at Spotify.
\textit{ACM RecSys}, s. 373.

\bibitem{brown2020}
Brown, T. et al., 2020.
Language Models are Few-Shot Learners.
\textit{NeurIPS}, 33, ss. 1877--1901.

\bibitem{apple2023}
Apple Inc., 2023.
MusicKit Framework Documentation.
\textit{Apple Developer Portal}.
Erişim Tarihi: 08 Haziran 2026.

\bibitem{fastapi}
Ramirez, S., 2021.
FastAPI: Modern, Fast Web Framework for Building APIs with Python.
\textit{Tiangolo Documentation}.
Erişim Tarihi: 02 Haziran 2026.

\end{thebibliography}

% ────────────────────────────────────────────────────────────
%  EKLER
% ────────────────────────────────────────────────────────────
\newpage
\phantomsection
\addcontentsline{toc}{chapter}{EKLER}
\chapter*{EKLER}

\section*{EK-1: Geliştirme Ortamı Gereksinimleri}

\begin{table}[H]
  \centering
  \caption*{Musaic Geliştirme Ortamı}
  \renewcommand{\arraystretch}{1.35}
  \begin{tabular}{ll}
    \hline
    \textbf{Bileşen}     & \textbf{Sürüm / Donanım} \\
    \hline
    Bilgisayar Donanımı  & MacBook Air (2020 Model) - M1 Çip \\
    İşletim Sistemi      & macOS Sonoma 14.0+ \\
    Geliştirme Ortamı    & Xcode 15+ (Swift 5.9) \\
    Hedef İşletim Sistemi& iOS 16.0+ \\
    Backend Ortamı       & Python 3.10+ (FastAPI) \\
    API Entegrasyonları  & Google Gemini, Firebase SDK, Apple MusicKit \\
    \hline
  \end{tabular}
\end{table}

\section*{EK-2: API Endpoint İstek Şemaları}

\begin{lstlisting}[language=Python, caption={Toplu Kaydırma (Batch Swipe) JSON İsteği}]
# /api/swipe - POST İstek Şeması
{
  "liked_songs": [
    {
      "song_id": 42,
      "title": "After Hours",
      "artist": "The Weeknd",
      "album_art": "https://...",
      "apple_music_id": "12345"
    }
  ],
  "disliked_song_ids": [10, 15, 23],
  "current_mood_vector": [0.5, 0.5, 0.4, 0.55, 0.3, 0.4, 0.1, 0.5, 0.1]
}

# /api/swipe - Sunucu Yanıtı
{
  "status": "ok",
  "new_vector": [0.48, 0.52, 0.38, 0.57, 0.28, 0.42, 0.08, 0.51, 0.09],
  "message": "Swipe session processed"
}
\end{lstlisting}

% ────────────────────────────────────────────────────────────
%  ÖZGEÇMİŞ
% ────────────────────────────────────────────────────────────
\newpage
\phantomsection
\addcontentsline{toc}{chapter}{ÖZGEÇMİŞ}
\chapter*{ÖZGEÇMİŞ}

\noindent\textbf{KİŞİSEL BİLGİLER}

\vspace{0.5cm}
\begin{tabular}{ll}
  Adı, Soyadı          & : Berke Balcı \\[4pt]
  Uyruğu               & : Türkiye (T.C.) \\[4pt]
  Doğum Tarihi ve Yeri & : -- / -- / ---- \\[4pt]
  Telefon              & : \\[4pt]
  Belgegeçer           & : \\[4pt]
  E-posta              & : 1030510273@erciyes.edu.tr \\[4pt]
  Adres                & : Melikgazi, KAYSERİ \\
\end{tabular}

\vspace{0.8cm}
\noindent\textbf{EĞİTİM}

\vspace{0.3cm}
\begin{tabular}{lll}
  \textbf{Derece} & \textbf{Kurum}                                 & \textbf{Mez. Yılı} \\[3pt]
  Lisans          & Erciyes Üniversitesi, Bilgisayar Mühendisliği  & 2026 \\[3pt]
  Lise            & ---                                            & \\
  Ortaokul        & ---                                            & \\
\end{tabular}

\vspace{1.5cm}
\noindent\rule{\linewidth}{0.5pt}
\vspace{0.5cm}

\noindent\textbf{KİŞİSEL BİLGİLER}

\vspace{0.5cm}
\begin{tabular}{ll}
  Adı, Soyadı          & : İlayda Baran \\[4pt]
  Uyruğu               & : Türkiye (T.C.) \\[4pt]
  Doğum Tarihi ve Yeri & : -- / -- / ---- \\[4pt]
  Telefon              & : \\[4pt]
  Belgegeçer           & : \\[4pt]
  E-posta              & : 1030510572@erciyes.edu.tr \\[4pt]
  Adres                & : Melikgazi, KAYSERİ \\
\end{tabular}

\vspace{0.8cm}
\noindent\textbf{EĞİTİM}

\vspace{0.3cm}
\begin{tabular}{lll}
  \textbf{Derece} & \textbf{Kurum}                                 & \textbf{Mez. Yılı} \\[3pt]
  Lisans          & Erciyes Üniversitesi, Bilgisayar Mühendisliği  & 2026 \\[3pt]
  Lise            & ---                                            & \\
  Ortaokul        & ---                                            & \\
\end{tabular}

\end{document}